\section*{Chapitre \ref{ch:m1:index-construction-and-rebalancing} --- Construction des indices et rééquilibrage}

\begin{solution}{exo:m1:index-construction-and-rebalancing:1}
Capitalisation totale de 20~milliards de dollars~; ajustée au flottant, $20 \times 0.65 = \$13$
poids $13/5\,000 = 0.26\%$.
\end{solution}

\begin{solution}{exo:m1:index-construction-and-rebalancing:2}
$\varphi = 2/5 = 40\%$ du flottant de chaque membre. Demande~: $0.40 \times 13 =
5.2$ milliards de dollars, soit treize jours de volume.
\end{solution}

\begin{solution}{exo:m1:index-construction-and-rebalancing:3}
Membre à 108~: il reste ($108 \le 110$). Valeur hors indice à 95~: elle n'entre pas (elle
a besoin de 90 ou mieux), sauf s'il reste des places à pourvoir par rang. Valeur hors indice à
88~: elle entre. Membre à 112~: il est exclu.
\end{solution}

\begin{solution}{exo:m1:index-construction-and-rebalancing:4}
Dans le grand indice~: la composition a été fixée le 30 avril, et les cours ultérieurs ne
modifient pas cette composition. La baisse modifie la \emph{taille} de l'ensemble~: le titre entre
avec un poids et une demande des trackers inférieure de 30~\% en dollars (le même nombre
d'actions, puisque les trackers détiennent une fraction $\varphi$ du flottant). La société
qui avait acheté en avril avait raison sur l'événement mais a perdu 30~\% sur le titre~: une
position sur indice est d'abord une position sur le titre. Le gain de quelques pour cent
lié à l'événement ne compense pas une exposition non couverte à la société.
\end{solution}

\begin{solution}{exo:m1:index-construction-and-rebalancing:5}
$10\,000 \times 1.5\% = 150$ milliards de dollars. Le membre à 6~\% devient $6\% \times
0.985 = 5.91\%$~: une vente de $0.09\% \times 10\,000 = 9$ milliards de dollars. Deux étapes
réduiraient de moitié la transaction à chaque clôture et la pression sur la liquidité~; mais chaque
étape est un événement distinct à anticiper, et l'indice porterait un
demi-poids qui ne représente rien pendant une semaine, et les trackers supporteraient
deux fois des décisions de réplication. Une seule étape sur la clôture la plus liquide
disponible concentre la liquidité là où se trouve la demande.
\end{solution}

\begin{solution}{exo:m1:index-construction-and-rebalancing:6}
Un jour~: $0.7 \times 2.5\% \times \sqrt{13} = 6.3\%$, de l'ordre des 7.4~\%
mesurés dans les années 1990 (et la règle est appliquée bien au-delà de son domaine).
Étalé sur vingt jours~: $0.7 \times 2.5\% \times \sqrt{13/20} = 1.4\%$ par jour
d'impact temporaire, du fait d'intervenants qui détiennent ensuite les actions à risque pendant
des semaines~; la concurrence entre eux et l'offre issue de l'indice que le titre
quitte, ramènent l'effet visible vers les 0.3~\% de la dernière décennie.
La même demande, absorbée par une file d'attente plus longue et plus encombrée.
\end{solution}

\begin{solution}{exo:m1:index-construction-and-rebalancing:7}
$b = 0$~: 18.5 ajouts par révision, taux de rotation à sens unique 1.92~\%. $b = 30$~: 9.2
ajouts, 1.18~\%. Le fonds négocie deux fois ce taux de rotation~; l'économie est
de $2 \times 0.74\% \times 15$ points de base $= 0.22$ point de base par an.
Les coûts explicites ne sont pas la raison des \omterm{def:m1:index-construction-and-rebalancing:buffer}{règles de tampon}~: les titres qui basculent à
la frontière sont les plus petits de l'indice. Ce que réduisent les tampons, c'est le nombre
d'événements prévisibles au cours desquels les trackers supportent un impact que les 15 points de base ne
capturent pas.
\end{solution}

\begin{solution}{exo:m1:index-construction-and-rebalancing:8}
(i) L'effet mesuré sur cette décennie est depuis retombé à quelques dixièmes de
pour cent~: la période de backtest est l'hypothèse. (ii) Le 5~\% a été mesuré
à des cours auxquels une petite étude pouvait négocier~; chaque intervenant achète désormais
à l'annonce, si bien que l'ouverture suivante intègre déjà le mouvement, et
vingt-cinq événements n'ont pas une capacité illimitée. (iii) Les rendements par événement
ne s'additionnent pas en un rendement annuel~: les événements se chevauchent, le capital est immobilisé pendant des jours à
des semaines, les positions portent un risque sur titre non couvert bien supérieur à l'avantage, et
les coûts ont été ignorés. On peut ajouter que de nombreux ajouts sont désormais anticipés
avant d'être annoncés, ce qui déplace le rendement hors de la fenêtre
de négociation retenue.
\end{solution}

\begin{solution}{pb:m1:index-construction-and-rebalancing:1}
\textbf{1.} $200/3\,000 = 6.67\%$; $4\,000/50\,000 = 8\%$.
\textbf{2.} $6.67\% \times 1.2$ milliard de dollars $= \$80$ millions~; 3.33 millions d'actions.
\textbf{3.} 3.7 jours.
\textbf{4.} Vente $6.67\% \times 9 = 600$ millions de dollars~; achat $8\% \times 9 =
720$ millions de dollars~; solde net $+120$ millions de dollars.
\textbf{5.} La vente représente 4.0 jours de volume, le solde net 0.8 jour. Si les deux groupes
négocient dans la même enchère de clôture, seul le solde net pèse sur le cours~: une
migration est surtout un croisement entre deux ensembles de trackers.
\textbf{6.} 2~millions d'actions contre une enchère normale de 72\,000~: 28 fois.
\textbf{7.} $0.7 \times 3\% \times \sqrt{3.7} = 4.0\%$.
\textbf{8.} $0.7 \times 3\% \times \sqrt{3.7/20} = 0.9\%$ par jour.
\textbf{9.} Les anticipateurs, qui ont acheté au cours des semaines précédentes
auprès de vendeurs ordinaires à un rythme que le marché pouvait absorber, et les \omterm{def:m1:what-a-trading-firm-does:market-maker}{teneurs de marché} qui
se mettent à découvert pendant l'enchère et rachètent dans les jours suivants. Le flux de
déséquilibre de l'enchère (voir au \cref{ch:m1:auctions}) leur indique combien il manque encore.
\textbf{10.} $0.8 \times 3\% - 0.2 \times 2\% - 0.4\% = 1.6\%$~: 80\,000~\$.
\textbf{11.} Le rendement vaut 3~\% ou $-2\%$~: écart-type $\sqrt{0.8
\times 0.2} \times 5\% = 2\%$, 100\,000~\$, avant tout risque ordinaire sur le titre.
\textbf{12.} 2~millions de dollars, 500\,000~\$, ratio 4.0.
\textbf{13.} Écart-type $100\,000 \times \sqrt{25 + 600 \times 0.3} =
1.43$ millions de dollars~; ratio 1.4.
\textbf{14.} $p = (0.4\% + 2\%)/5\% = 0.48$.
\textbf{15.} Toutes les positions sont longues sur de petites valeurs proches d'un même seuil~: une
variation de marché avant la date de classement déplace le seuil et toutes les
probabilités ensemble~; et les mêmes anticipateurs détiennent tous les titres, si bien que
lorsque l'un d'eux réduit son risque, chaque position baisse en même temps.
\textbf{16.} Chaque événement touche moins de titres, avec des poids plus faibles, puisque la
dérive accumulée entre deux révisions est moitié moindre~; il y a deux fois plus de dates.
Le volume total à négocier baisse un peu (moins de titres s'écartent fortement de leur indice),
et l'activité devient moins saisonnière et moins concentrée sur un seul vendredi.
\textbf{17.} On les juge sur le \omterm{def:m1:the-buy-side:benchmark}{tracking error} par rapport à un indice calculé à
cette clôture. Acheter tôt à meilleur prix constitue un gain de quelques points de base
s'il fonctionne, et un écart inexplicable s'il ne fonctionne pas.
\textbf{18.} Les investisseurs du fonds, à travers le \omterm{def:m1:the-buy-side:benchmark}{tracking error} dans les deux
sens~; ils le découvrent par l'écart de suivi dans le rapport annuel,
et de nulle autre source.
\textbf{19.} \textbf{3.7 jours de volume moyen.}
\textbf{20.} Ils jouent le rôle de réservoir qui transforme une demande d'une minute en une
demande de trois semaines, et leur profit en est la rémunération.
\end{solution}

\begin{solution}{iq:m1:index-construction-and-rebalancing:1}
Les trackers doivent acheter une fraction connue de son flottant à la clôture effective.
Historiquement, le titre bondissait à l'annonce et montait jusqu'à la
date d'effet, avec un repli partiel. Pour le grand indice américain, l'effet mesuré
à l'annonce est passé de plusieurs pourcents dans les années 1990 à presque
rien~: le mouvement se produit désormais plus tôt, car les ajouts sont anticipés, et la
demande est satisfaite par les firmes qui se positionnent à l'avance et par les trackers
de l'indice que le titre quitte.
\iqlookfor{la chronologie (prédiction, annonce, clôture effective), et
la connaissance que l'effet de manuel s'est réduit.}
\end{solution}

\begin{solution}{iq:m1:index-construction-and-rebalancing:2}
Un indice est censé pouvoir être détenu par tous simultanément. Avec une pondération
intégrale par capitalisation, les trackers auraient besoin d'actions qu'un fondateur ou un État
ne vendra jamais~: pour une société détenue à 80~\% par son fondateur, une demande
cinq fois plus forte rapportée aux actions négociables que celle de ses voisines, un squeeze à
chaque entrée de flux, et un cours déformé par l'indice lui-même.
\iqlookfor{l'argument du «~tous en même temps~» et une compression concrète.}
\end{solution}

\begin{solution}{iq:m1:index-construction-and-rebalancing:3}
$8\% \times 10$ milliards $= 800$ millions de dollars. Savoir si c'est beaucoup dépend du
volume quotidien du titre (un jour ou quinze), de la taille de sa
enchère de clôture, du fait qu'il quitte un autre indice doté de ses propres trackers (le
solde net, et non le brut, compte), et de la part déjà achetée par les anticipateurs, que
révéleront la hausse préalable du cours et l'encours vendeur à découvert de
l'enchère.
\iqlookfor{le raccourci par fraction du flottant et le solde net de la migration.}
\end{solution}

\begin{solution}{iq:m1:index-construction-and-rebalancing:4}
Reconstruire l'univers du fournisseur et ses capitalisations totales à partir de la
méthodologie~; classer~; localiser les points de rupture et appliquer les bandes de tampon aux
membres actuels~; avant la date de classement, associer à chaque valeur frontière une
probabilité de franchissement, fondée sur la distance et la volatilité. Erreurs~: nombre d'actions (classes multiples,
émissions récentes, déclarations que le fournisseur interprète autrement), détails d'éligibilité
(domiciliation, flottant minimal, cotation), les points de rupture eux-mêmes qui évoluent avec
le marché, et les opérations sur titres entre la date de classement et la date d'effet.
\iqlookfor{la prudence sur le nombre d'actions et le point de rupture mobile.}
\end{solution}

\begin{solution}{iq:m1:index-construction-and-rebalancing:5}
L'effet mesuré est le rendement sur une fenêtre qui commence à
l'annonce. Davantage d'actifs passifs signifie une demande plus forte, mais aussi une
rémunération plus élevée pour l'anticiper et la fournir~: le cours s'ajuste avant la
fenêtre, davantage de capitaux se disputent le rôle de vendeur à la clôture, et de nombreux ajouts
migrent depuis un indice sœur dont les trackers sont des vendeurs naturels. Le choc
a augmenté~; son prix a baissé.
\iqlookfor{la distinction entre la demande et le rendement mesuré.}
\end{solution}

\begin{solution}{iq:m1:index-construction-and-rebalancing:6}
Risques~: erreur de prévision (le titre n'est pas ajouté)~; risque sur le titre entre l'entrée et
la clôture, qui dépasse largement l'avantage et doit être couvert par secteur et par taille~;
le crowding, visible sous forme de baisses corrélées de tous les candidats lorsqu'un
concurrent réduit son risque~; points de rupture qui évoluent avec le marché~; modifications de règles par le
fournisseur~; et l'exécution dans une seule enchère. Dimensionnement~: en jours de volume
pour sortir, et non en dollars~; par l'avantage pondéré par la probabilité, net du coût de couverture~; avec un
plafond sur la part de la demande attendue des trackers que le livre peut détenir, puisque
au-delà d'une certaine part, le livre devient sa propre sortie.
\iqlookfor{la couverture de la jambe action, le crowding et le dimensionnement
en unités de liquidité.}
\end{solution}
